\section{Conclusion}
\label{sec:conclusion}

This work separated multi-person pose grouping from keypoint detection and measured what the separation costs.
PG-GAT --- a graph attention network with a type-pair attention bias, skeleton-relative position encoding,
and same-type repulsion heads --- embeds the detected keypoints so that a plain $k$-means read-out recovers the people.
Trained on synthetic scenes and fine-tuned on COCO,
it reaches a pose grouping accuracy of 0.9715 on ground-truth keypoints and 0.9591 on HigherHRNet detections,
against 0.9847 for the detector's own jointly trained grouping.
With a K-head trained on the frozen embeddings, the pipeline runs with no ground truth at inference,
and the same frozen network, applied zero-shot to YOLO11x-pose and RTMO-L, trails each detector's native grouping by about 0.02.
The cost is larger at the system level, 0.048 AP on identical joints, and it concentrates in crowded scenes.
Across the study, changes to the embedding moved accuracy far more than changes to the grouping stage:
learned heads and modularity-based clustering did not improve on $k$-means, while the skeleton-aware attention,
the architecture it was sized to, and fine-tuning on real data did.

The results point to three directions.
Extending the synthetic data beyond five people addresses the largest measured deficit,
and representing structure above the joint, so that a wrist is disambiguated through its limb,
targets the errors that concentrate at the extremities.
Repeating the key comparisons over multiple seeds would establish
which of the smaller differences reported here are systematic.